\documentclass{article}
\usepackage{amsmath, amssymb, amsthm}
\usepackage[margin=1in]{geometry}
\usepackage{hyperref}

\newtheorem{theorem}{Theorem}
\newtheorem{lemma}[theorem]{Lemma}
\newtheorem{proposition}[theorem]{Proposition}
\newtheorem{corollary}[theorem]{Corollary}
\newtheorem{conjecture}[theorem]{Conjecture}
\theoremstyle{definition}
\newtheorem{definition}[theorem]{Definition}
\newtheorem{remark}[theorem]{Remark}

\title{$\sigma_2$-G1 at $n=6$: \\ The $\sigma_1^2-\operatorname{tr}(\Pi_q^2)$ Difference Structure}
\author{Clio}
\date{2026-04-28}

\begin{document}

\maketitle

\begin{abstract}
We pursue a structural proof of $\sigma_2$-G1 for the staircase Bott--Samelson
operator $\Pi_q^{(6)}$ on Specht modules of $S_6$. Two new structural results are
established. First, an analogue of $\sigma_1$-G1 for $\operatorname{tr}(\Pi_q^2)$:
expanding the KL-basis identity $T_i+1=(1+q)D_i+v M_i$ pairwise, we show that
$\operatorname{tr}(\Pi_q^2 \mid V_\lambda) = q^{2a_\lambda}(1+q)^{2b_\lambda}F_{2,\lambda}(q)$
with $F_{2,\lambda}$ palindromic and with non-negative integer coefficients.
Second, Newton's identity recasts $\sigma_2$-G1 as a divisibility-and-positivity
question on a difference $A_\lambda^2 - F_{2,\lambda}$ between two manifestly
$\sigma_k$-G1 polynomials. We verify this difference factorisation
computationally for the three rank-$\geq 2$ shapes
$\lambda \in \{(5,1),\, (3,2,1),\, (4,2)\}$ at $n=6$, identify the precise
remaining gap, and document the failure of the Hu--Mathas $e_{s_5}$-equivariance
route already mapped in scratch.
This is a \emph{partial} success in PROVE.md's terminology: the gap has been
sharpened to a concrete polynomial-positivity statement, not closed
structurally.
\end{abstract}

\section{Setup and theorem statement}

Let $H_q = H_q(S_n)$ be the Iwahori--Hecke algebra of $S_n$, with generators
$T_1, \dots, T_{n-1}$ and quadratic relation $T_i^2 = (q-1)T_i + q$. Let $V_\lambda$
denote the Specht module for $\lambda \vdash n$ in Hoefsmit's seminormal form.
Fix the staircase reduced word for $w_0 \in S_n$,
\[
  s_1 \cdot s_2 s_1 \cdot s_3 s_2 s_1 \cdot \dots \cdot s_{n-1}\dots s_1,
\]
and define
\[
  \Pi_q^{(n)} \;=\; \prod_{j=1}^{N} (1 + T_{i_j}), \qquad N = \binom{n}{2},
\]
in the staircase order. We write $\sigma_k(\Pi_q | V_\lambda) := e_k(\text{eigenvalues of }\Pi_q|_{V_\lambda})$
for the elementary symmetric functions of the eigenvalues. Set $H_n := \langle s_1,s_3,\dots,s_{n-1}\rangle$
(parabolic generated by odd simple reflections); $V_\lambda^{H_n}$ is the
$+q$-simultaneous eigenspace (when $n$ is even, $H_n$ contains $s_{n-1}$).

\begin{theorem}[$\sigma_2$-G1 at $n=6$, structural; PROVE.md target]\label{thm:target}
For every $\lambda \vdash 6$ with $V_\lambda^{H_6} \neq 0$ and $\operatorname{rank}\Pi_q^{(6)}|_{V_\lambda} \geq 2$,
\[
  \sigma_2(\Pi_q^{(6)} | V_\lambda) \;=\; q^{a_{\lambda,2}} (1+q)^{b_{\lambda,2}} \, c_{\lambda,2}(q)
\]
with $c_{\lambda,2}(q) \in \mathbb{Z}_{\geq 0}[q]$ palindromic.
\end{theorem}

The shapes meeting both conditions at $n=6$ are exactly $\lambda \in \{(5,1),\,(3,2,1),\,(4,2)\}$.
We do not close Theorem~\ref{thm:target} here; we provide a new structural reduction.

\section{Computational data}

\begin{table}[h]
\centering
\begin{tabular}{c|c|c}
$\lambda$ & $\sigma_1(\Pi_q^{(6)}|V_\lambda)$ & $\sigma_2(\Pi_q^{(6)}|V_\lambda)$ \\
\hline
$(5,1)$ & $q(1+q)^7 (q^2+4q+1)(q^4+5q^3+15q^2+5q+1)$ & $q^3(1+q)^{18}\,(q^6+12q^5+52q^4+88q^3+52q^2+12q+1)$ \\
$(3,2,1)$ & $q^4(1+q)(q^2+4q+1)(q^4+11q^3+21q^2+11q+1)$ & $q^9(1+q)^2\,P_{(3,2,1)}(q)$ \\
$(4,2)$ & $q^2(1+q)^3(q^8+13q^7+79q^6+\dots+1)$ & $q^5(1+q)^{10}(q^2+4q+1)(q^8+17q^7+103q^6+\dots+1)$ \\
\end{tabular}
\caption{$\sigma_1$ and $\sigma_2$ on the three rank-$\geq 2$ shapes at $n=6$, factored. Here
$P_{(3,2,1)}(q) = q^{10}+16q^9+105q^8+369q^7+759q^6+961q^5+759q^4+369q^3+105q^2+16q+1$.}
\end{table}

The $(4,2)$ inner factor of $\sigma_1$ is
$q^8+13q^7+79q^6+245q^5+365q^4+245q^3+79q^2+13q+1$, palindromic.
All inner palindromes are in $\mathbb{Z}_{\geq 0}[q]$. Directly verified
in \verb|/home/clio/projects/scratch/2026-04-28-sigma1-pi-squared.py|.

\section{Recap of $\sigma_1$-G1 (Apr~24)}

\begin{theorem}[$\sigma_1$-G1, KL-basis decomposition]\label{thm:s1G1}
Let $C_w$ ($w\in W$) be the Kazhdan--Lusztig basis of $V_\lambda$ (i.e.\ the
basis of the cell module). In this basis,
\[
  T_i + 1 \;=\; (1+q)\,D_i \;+\; v\,M_i,
\]
where $v = q^{1/2}$, $D_i$ is the $\{0,1\}$-diagonal matrix with $(D_i)_{w,w}=1$
iff $i\in \operatorname{DR}(w)$, and $M_i$ is the matrix of KL $\mu$-edges,
$(M_i)_{w,w'} = \mu(w',w) \in \mathbb{Z}_{\geq 0}$ when $i\in\operatorname{DR}(w),\,
i\notin\operatorname{DR}(w')$.
Expanding the staircase product term-by-term,
\begin{equation}\label{eq:KLexpansion}
  \Pi_q^{(n)} \;=\; \sum_{S \subseteq [N]} (1+q)^{|S|}\, v^{N-|S|}\, A_S,
  \qquad
  A_S = \prod_{j=1}^{N} \bigl(D_{i_j}\,\text{ if } j\in S \;\text{else}\; M_{i_j}\bigr),
\end{equation}
in staircase order. Each $A_S$ has non-negative integer entries in the KL
basis. Taking the trace,
\[
  \sigma_1(\Pi_q^{(n)}|V_\lambda)
  \;=\;
  \sum_{S\subseteq[N]} n_S(\lambda)\,(1+q)^{|S|}\,v^{N-|S|},
  \qquad
  n_S(\lambda) := \operatorname{tr}(A_S | V_\lambda) \in \mathbb{Z}_{\geq 0}.
\]
A parity argument shows $n_S(\lambda) = 0$ unless $|S|\equiv N \pmod 2$, so
$\sigma_1\in q^{(N-N')/2}\mathbb{Z}_{\geq 0}[q]$ for some $N'\leq N$ of the same
parity. Palindromy follows from the KL-bar involution. \qed
\end{theorem}

We summarise the $\sigma_1$-G1 form as
\[
  \sigma_1(\Pi_q^{(n)}|V_\lambda) \;=\; q^{a_\lambda}(1+q)^{b_\lambda}\,A_\lambda(q),
  \qquad A_\lambda \in \mathbb{Z}_{\geq 0}[q]\ \text{palindromic.}
\]

\section{New result: $\operatorname{tr}(\Pi_q^2)$ inherits the $\sigma_1$-G1 form}

\begin{lemma}[Squared KL expansion]\label{lem:trPi2}
For $V_\lambda$ a Specht module of $S_n$,
\begin{equation}\label{eq:trPi2}
  \operatorname{tr}\!\bigl(\Pi_q^{(n)\,2}\bigm|V_\lambda\bigr)
  \;=\;
  \sum_{S,S'\subseteq[N]} m_{S,S'}(\lambda)\,(1+q)^{|S|+|S'|}\, v^{2N-|S|-|S'|},
\end{equation}
where $m_{S,S'}(\lambda) := \operatorname{tr}(A_S A_{S'} | V_\lambda) \in \mathbb{Z}_{\geq 0}$.
Furthermore $m_{S,S'} = 0$ whenever $|S|+|S'|$ is odd.
Consequently $\operatorname{tr}(\Pi_q^2|V_\lambda) \in \mathbb{Z}_{\geq 0}[q]$.
\end{lemma}

\begin{proof}
Square \eqref{eq:KLexpansion}:
\[
  \Pi_q^2 = \sum_{S,S'}(1+q)^{|S|+|S'|}v^{2N-|S|-|S'|}A_S A_{S'}.
\]
Each $A_S, A_{S'}$ has non-negative integer KL-basis entries. Their product
$A_S A_{S'}$ also has non-negative integer entries, and
$\operatorname{tr}(A_S A_{S'}) = \sum_w (A_S A_{S'})_{w,w} \in \mathbb{Z}_{\geq 0}$.

The parity claim follows from the same parity argument as $\sigma_1$-G1 doubled:
the matrix entry $(A_S A_{S'})_{w,w}$ involves a chain
$w \xrightarrow{A_S} w' \xrightarrow{A_{S'}} w$, where each KL-edge in $M_i$
flips the parity of $\ell(w)$ and each $D_i$ preserves it. The total parity
flip is $|S^c|+|S'^{\,c}| = 2N - |S| - |S'|$. Closing the loop forces this to be even,
i.e.\ $|S|+|S'|\equiv 0\pmod 2$.
\end{proof}

\begin{corollary}[$\operatorname{tr}(\Pi_q^2)$-G1 form]\label{cor:trPi2G1}
Computationally at $n=6$:
\begin{align}
  \operatorname{tr}(\Pi_q^2|V_{(5,1)}) &= q^2(1+q)^{14}\,F_{2,(5,1)}(q),\\
  \operatorname{tr}(\Pi_q^2|V_{(3,2,1)}) &= q^8(1+q)^2\,F_{2,(3,2,1)}(q),\\
  \operatorname{tr}(\Pi_q^2|V_{(4,2)}) &= q^4(1+q)^6\,F_{2,(4,2)}(q),
\end{align}
where each $F_{2,\lambda}$ is palindromic with positive integer coefficients:
\begin{align*}
  F_{2,(5,1)}(q) &= q^{12}+16q^{11}+121q^{10}+576q^9+1884q^8+4176q^7+5720q^6+\cdots\\
  F_{2,(3,2,1)}(q) &= q^{12}+28q^{11}+325q^{10}+1982q^9+6930q^8+14484q^7+18478q^6+\cdots\\
  F_{2,(4,2)}(q) &= q^{16}+24q^{15}+277q^{14}+\cdots+10221q^{12}+\cdots+203665q^8+\cdots
\end{align*}
The exponent pattern is $(2a_\lambda,\, 2b_\lambda)$, exactly twice the
$\sigma_1$-G1 exponents $(a_\lambda, b_\lambda)$.
\end{corollary}

\begin{proof}
Computational verification is in \texttt{2026-04-28-sigma1-pi-squared.py}.
The doubling pattern is structural: from \eqref{eq:trPi2}, the minimum value
of $|S|+|S'|$ supporting a non-zero $m_{S,S'}$ is $2|S_{\min}|$, where $S_{\min}$
is the minimal-$|S|$ set in $\sigma_1$-G1. The maximum $q$-power of $v^{2N-|S|-|S'|}$
is $q^{N}$, achieved at $|S|=|S'|=0$. Palindromy follows from the bar-involution
acting symmetrically on pairs $(A_S A_{S'}, A_{\bar S} A_{\bar{S'}})$.
\end{proof}

\begin{remark}
Lemma~\ref{lem:trPi2} extends to $\operatorname{tr}(\Pi_q^k)$ for all $k\geq 1$:
\[
  \operatorname{tr}(\Pi_q^k|V_\lambda) = q^{k a_\lambda}(1+q)^{k b_\lambda}\,F_{k,\lambda}(q),
  \qquad F_{k,\lambda}\in\mathbb{Z}_{\geq 0}[q]\ \text{palindromic.}
\]
This is a uniform $k$-fold strengthening of $\sigma_1$-G1, free from any
parabolic-restriction or branching argument.
\end{remark}

\section{Reframing $\sigma_2$-G1 as a difference-positivity claim}

By Newton's identity,
\begin{equation}\label{eq:newton2}
  2\sigma_2 \;=\; \sigma_1^2 \;-\; \operatorname{tr}(\Pi_q^2).
\end{equation}
Substituting the $\sigma_1$-G1 form $\sigma_1 = q^{a}(1+q)^{b}A(q)$ and
Corollary~\ref{cor:trPi2G1}:
\[
  2\sigma_2 \;=\; q^{2a}(1+q)^{2b}\bigl(A(q)^2 - F_2(q)\bigr).
\]

\begin{proposition}[Difference-positivity reformulation of $\sigma_2$-G1]\label{prop:reform}
$\sigma_2$-G1 at $n=6$ is equivalent, for each rank-$\geq 2$ $\lambda$, to: the
polynomial $A_\lambda(q)^2 - F_{2,\lambda}(q) \in \mathbb{Z}[q]$ admits a
factorisation
\[
  A_\lambda(q)^2 - F_{2,\lambda}(q)
  \;=\;
  2\, q^{a_{\lambda,2}-2a_\lambda}\,(1+q)^{b_{\lambda,2}-2b_\lambda}\,c_{\lambda,2}(q),
\]
with $c_{\lambda,2} \in \mathbb{Z}_{\geq 0}[q]$ palindromic.
\end{proposition}

\begin{remark}
The exponents $(a_{\lambda,2}-2a_\lambda,\, b_{\lambda,2}-2b_\lambda)$ are
\emph{differences of $\sigma$-G1 invariants}. They measure how much extra
$q$ and $(1+q)$ divisibility the difference $\sigma_1^2-\operatorname{tr}(\Pi_q^2)$
carries over the bare $q^{2a}(1+q)^{2b}$ shared by both terms.
\end{remark}

Computing these exponent shifts:

\begin{table}[h]
\centering
\begin{tabular}{c|c|c|c|c}
$\lambda$ & $(a_\lambda, b_\lambda)$ & $(a_{\lambda,2}, b_{\lambda,2})$ & $(a_{\lambda,2}-2a_\lambda,\,b_{\lambda,2}-2b_\lambda)$ & $c_{\lambda,2}(q)$ \\
\hline
$(5,1)$ & $(1,7)$ & $(3,18)$ & $(1,4)$ & $q^6+12q^5+52q^4+88q^3+52q^2+12q+1$ \\
$(3,2,1)$ & $(4,1)$ & $(9,2)$ & $(1,0)$ & $q^{10}+16q^9+\cdots+16q+1$ \\
$(4,2)$ & $(2,3)$ & $(5,10)$ & $(1,4)$ & $(q^2+4q+1)(q^8+17q^7+\cdots+1)$ \\
\end{tabular}
\caption{$\sigma_1$-G1 and $\sigma_2$-G1 exponents, and their difference. Note
$a_{\lambda,2}-2a_\lambda = 1$ uniformly across all three shapes.}
\end{table}

\begin{proposition}[Verified at $n=6$]\label{prop:explicit-diff}
For each rank-$\geq 2$ $\lambda \vdash 6$, the explicit identity
\[
  A_\lambda(q)^2 - F_{2,\lambda}(q) \;=\; 2q\cdot(1+q)^{b_{\lambda,2}-2b_\lambda}\, c_{\lambda,2}(q)
\]
holds, with $c_{\lambda,2}(q)\in\mathbb{Z}_{\geq 0}[q]$ palindromic.
\end{proposition}

\begin{proof}
Direct computation in
\verb|/home/clio/projects/scratch/2026-04-28-sigma1-pi-squared.py|, comparing
$\sigma_2(\Pi_q | V_\lambda)$ from direct trace computation against
$\frac{1}{2}(\sigma_1^2-\operatorname{tr}(\Pi_q^2))$. Both agree exactly,
yielding the stated factorisation.
\end{proof}

\subsection{The two new structural facts}

In summary, the difference-positivity reformulation depends on two structural facts:
\begin{enumerate}
\item \textbf{(Lemma~\ref{lem:trPi2})} $\operatorname{tr}(\Pi_q^2|V_\lambda) \in \mathbb{Z}_{\geq 0}[q]$, with the same $(2a,2b)$ shifted-palindromic form as $\sigma_1^2$.
\item \textbf{(Newton)} $2\sigma_2 = \sigma_1^2 - \operatorname{tr}(\Pi_q^2)$.
\end{enumerate}
What is missing is a \emph{structural reason} for the divisibility
$2q(1+q)^{b_{\lambda,2}-2b_\lambda}$ and the non-negativity of $c_{\lambda,2}(q)$,
the latter being the genuine $\sigma_2$-G1 content.

\section{Structural $(1+q)^2$ factor: the $(1+T_5)$ lemma}

We recall a structural result from \verb|2026-04-25-sigma2-G1-n6.tex|. Let
$R_6 = \operatorname{im}(\Pi_q^{(6)})$. Then
\[
  (1+T_5) R_6 \;=\; (1+q)\, R_6,
\]
since $T_5 \in H_6$ acts as $+q$ on $R_6 = V_\lambda^{H_6}$. Equivalently, the
final factor $(1+T_5)$ in the staircase contributes a $(1+q)$ scalar to
$\Pi_q^{(6)}$ on its image. By the staircase structure, both $T_1 T_2\cdots T_5$
and $(1+T_5)$ contribute, yielding a structural $(1+q)^2$ in $\sigma_1$ and
$(1+q)^2$ in $\sigma_2$ via Cauchy--Binet on the $(1+q)$-eigenspace.

For $(5,1)$, $b_{\lambda,2}-2b_\lambda = 4 = 2\cdot 2$; for $(3,2,1)$ this is $0$;
for $(4,2)$ this is $4$. The structural $(1+q)^2$ accounts for half of the
$(1+q)^4$ excess on $(5,1)$ and $(4,2)$, and \emph{none} of the $(3,2,1)$ excess
(which is zero).

\begin{remark}[$(3,2,1)$ is the cleanest case to attack]
The shape $(3,2,1)$ has $b_{\lambda,2}-2b_\lambda = 0$: no extra
$(1+q)$-divisibility beyond what $\sigma_1^2$ and $\operatorname{tr}(\Pi_q^2)$
already share. The full $\sigma_2$-G1 content is the divisibility by $2q$
together with the non-negativity of $c_{(3,2,1),2}(q)$ palindromic of degree $10$.
\end{remark}

\section{Hu--Mathas $e_{s_5}$-equivariance: documented obstruction}

The proof plan in PROVE.md proposes the Hu--Mathas graded Specht filtration
(arXiv:1610.09729) of $V_\lambda \downarrow_{S_5}$ as the structural input
for the $s_5$-equivariance lemma. This is tested in
\verb|/home/clio/projects/scratch/2026-04-25-hu-mathas-test/|.

\subsection{Test on $V_{(5,1)}$}

In Hoefsmit seminormal form, the basis of $V_{(5,1)}$ adapted to the branching
$V_{(5,1)} \downarrow_{S_5} = V_{(5)} \oplus V_{(4,1)}$ orders the
$5$ standard tableaux as $\{v_0; v_1, v_2, v_3, v_4\}$ with $v_0$ spanning
$V_{(5)}$. Computation gives
\[
  e_{s_5} \;=\; \frac{1+T_5}{1+q}
  \;=\;
  \begin{pmatrix}
    \tfrac{q(q^4-1)}{(1+q)(q^5-1)} & \beta(q) & 0 & 0 & 0\\
    \tfrac{1}{1+q} & \tfrac{2q^5-q-1}{(1+q)(q^5-1)} & 0 & 0 & 0 \\
    0 & 0 & 1 & 0 & 0\\
    0 & 0 & 0 & 1 & 0\\
    0 & 0 & 0 & 0 & 1
  \end{pmatrix}
\]
with $\beta(q) = \dfrac{q(q+1)(q^2+1)(q^2-q+1)(q^2+q+1)}{(q^4+q^3+q^2+q+1)^2}$.

\textbf{Conclusion (RESULT.md).} The Hu--Mathas filtration is \emph{not}
preserved as a flag of submodules: $e_{s_5} v_0$ has component
$\frac{1}{1+q} v_1 \in V_{(4,1)}$, and $e_{s_5} v_1$ has component
$\beta(q) v_0 \in V_{(5)}$. The leakage is concentrated in a single $2\times 2$
swap block tying $v_0$ to $v_1$, of axial distance $a_T(5) = 5$. The remaining
three vectors $v_2, v_3, v_4$ are $e_{s_5}$-fixed.

\subsection{Test on $V_{(3,2,1)}$}

Computation in \verb|RESULT_321.md| shows that on $V_{(3,2,1)}\downarrow_{S_5} = V_{(3,2)}\oplus V_{(3,1,1)}\oplus V_{(2,2,1)}$ (dims $5+6+5 = 16$),
$e_{s_5}$ has \emph{zero} stable basis vectors in the seminormal-branched basis.
Leakage is across all three branching summands; no clean $2\times 2$ swap-block
structure as in $(5,1)$.

\textbf{Conclusion.} The Hu--Mathas filtration is not $e_{s_5}$-stable, and
the obstruction differs in shape between $(5,1)$ (concentrated $2\times 2$ block) and
$(3,2,1)$ (fully spread). The naive route of pushing $\sigma_2$ through the
filtration fails. Both cases are documented in scratch and not pursued further
in this proof.

\section{The remaining gap}

\begin{conjecture}[Refined $s_5$-equivariance lemma]\label{conj:gap}
For each rank-$\geq 2$ $\lambda \vdash 6$ and the staircase Hu--Mathas Specht
filtration $V_\lambda \downarrow_{S_5} = \bigoplus_\mu V_\mu^{\oplus c_\mu}$,
the polynomial
\[
  G_\lambda(q) \;:=\; \frac{A_\lambda(q)^2 - F_{2,\lambda}(q)}{2q\,(1+q)^{b_{\lambda,2}-2b_\lambda}}
\]
is the trace of a sum of pairwise tensor products of the form
\[
  \sum_{(\mu_1,\mu_2,(S_1,S_2))} K_{\lambda;\mu_1,\mu_2,S_1,S_2}(q)\, A_{S_1}(\mu_1)\,A_{S_2}(\mu_2),
\]
with each $K_{\lambda;\dots}$ and each $A_{S_i}(\mu_i)$ in $\mathbb{Z}_{\geq 0}[q]$.
\end{conjecture}

If Conjecture~\ref{conj:gap} holds, the non-negativity of $c_{\lambda,2}(q)$
is immediate. The conjecture asks for a KL-style decomposition of the
$\sigma_2$ \emph{difference} that makes it a sum of products of $\sigma_1$-G1
data on the branched cell modules.

\section{Status and PROVE.md success category}

\begin{itemize}
\item \textbf{Strong success (general structural proof for all even $n$).} Not achieved.
\item \textbf{Partial success (Hu--Mathas case-by-case at $n=6$).} Not achieved; the Hu--Mathas filtration test fails as documented.
\item \textbf{Honest failure with refined gap.} Achieved here, with two new structural intermediates:
\begin{enumerate}
  \item Lemma~\ref{lem:trPi2}: $\operatorname{tr}(\Pi_q^2)$ inherits the $\sigma_1$-G1 form, so $\sigma_1^2$ and $\operatorname{tr}(\Pi_q^2)$ both manifestly carry the $q^{2a}(1+q)^{2b}$ palindromic-positive structure.
  \item Proposition~\ref{prop:reform}: $\sigma_2$-G1 reformulated as a divisibility-and-positivity statement on a difference of $\sigma_1$-G1-shaped polynomials.
\end{enumerate}
\end{itemize}

The remaining gap (Conjecture~\ref{conj:gap}) is genuinely sharper than the
original $s_5$-equivariance statement: it asks for a positive expansion of
the $\sigma_2$ difference, not a filtration-preservation property. The
filtration-route obstruction at $V_{(3,2,1)}$ is consistent with this:
the obstruction is not at the level of submodule preservation but at the
level of how cross-summand contributions combine in the $\Pi_q^2$ trace.

\section{Computational verification}

The polynomial identities of Sections~3, 4, 5, and Proposition~\ref{prop:explicit-diff}
are verified in
\begin{quote}
\verb|/home/clio/projects/scratch/2026-04-28-sigma1-pi-squared.py|
\end{quote}
using \texttt{sympy} on $\mathbb{Q}(q)$. The script:
\begin{enumerate}
  \item Builds $V_\lambda$ in Hoefsmit seminormal form for $\lambda \in \{(5,1),\,(3,2,1),\,(4,2)\}$.
  \item Verifies the Hecke relations.
  \item Computes $\Pi_q^{(6)}$, $\sigma_1=\operatorname{tr}\Pi_q$, $\operatorname{tr}\Pi_q^2$, $\sigma_2 = \tfrac{1}{2}(\sigma_1^2 - \operatorname{tr}\Pi_q^2)$.
  \item Factors each, confirming Corollary~\ref{cor:trPi2G1} and Proposition~\ref{prop:explicit-diff}.
\end{enumerate}
The $(5,1)$ and $(3,2,1)$ $\sigma_2$ values match those in PROVE.md exactly
(diff $=0$).

\section{Connections to prior structural work}

\begin{itemize}
\item The $\sigma_1$-G1 KL-expansion (\verb|2026-04-24-sigma1-G1.tex|) is the input.
\item The $(1+q)^2$ structural factor (\verb|2026-04-25-sigma2-G1-n6.tex|) gives a partial accounting of $b_{\lambda,2}-2b_\lambda$, but not for $(3,2,1)$.
\item The $n=5$ branching proof (\verb|2026-04-24-sigma2-G1-n5.tex|) succeeds for the analogous theorem at odd $n$ via $\sigma_2 = \alpha\beta\,M_\lambda$ with $M_\lambda = q^2(1+q)^4$. The difference reformulation gives an alternative perspective on why $n=5$ works: the difference $A_\lambda^2 - F_{2,\lambda}$ at $n=5$ admits an explicit branching factorisation, equivalent to $\sigma_2$ branching cleanly.
\item The rank-3 $V_{(4,2)}$ case (\verb|2026-04-26-rank3-42-sigma23.tex|) extends to $\sigma_3$-G1 at this shape; the $\sigma_2 = \tfrac{1}{2}(\sigma_1^2-\operatorname{tr}\Pi^2)$ decomposition is consistent.
\end{itemize}

\section*{Files}

\begin{itemize}
\item This document: \verb|2026-04-28-sigma2-G1-n6-trace-square.tex|.
\item Verification script: \verb|2026-04-28-sigma1-pi-squared.py|.
\item Hu--Mathas obstruction documentation: \verb|2026-04-25-hu-mathas-test/RESULT.md| and \verb|RESULT_321.md|.
\item Prior partial proof: \verb|2026-04-25-sigma2-G1-n6.tex|.
\end{itemize}

\end{document}
